\documentclass[conference]{IEEEtran}

\usepackage{cite}
\usepackage{amsmath,amssymb,amsfonts}
\usepackage{algorithmic}
\usepackage{graphicx}
\usepackage{textcomp}
% extra colors used with minted
\usepackage[svgnames]{xcolor}
\usepackage{booktabs}
\usepackage{tagging}
% for code blocks
\usepackage{minted}
% to format description
\usepackage{enumitem}

% --- environments for verbatim blocks ---
\newenvironment{bash}
{\minted[
frame=lines,
framesep=2mm,
baselinestretch=0.9,
bgcolor=WhiteSmoke,
fontsize=\footnotesize]{bash}}
{\endminted}

\newenvironment{example}
{\minted[bgcolor=GhostWhite,fontsize=\small]{text}}
{\endminted}

\usepackage{sc24repro}


% Uncomment/comment the following usetag lines 
% if you would like to get an explanation or an example.
% Don't forget to comment them for the final submission.
% \usetag{explanation}
% \usetag{example}

\begin{document}

\twocolumn[%
    {\begin{center}
                \Huge
                Appendix: Artifact Description/Artifact Evaluation
            \end{center}}
]

%%%%%%%%%%%%%%%%%%%%%%%%%%%%%%%%%%%%%%%%%%%%%%%%%%%%%
%  AD Appendix
%%%%%%%%%%%%%%%%%%%%%%%%%%%%%%%%%%%%%%%%%%%%%%%%%%%%%

\appendixAD

\section{Overview of Contributions and Artifacts}

\subsection{Paper's Main Contributions}

\begin{description}
    \item{$C_1$} CARP: an in-situ indexing system that builds range-query indexes for scientific application output data as the data is written to storage.
    \item{$C_2$} An analysis of key distributions generated by the output of the VPIC scientific simulation.
\end{description}


\subsection{Computational Artifacts}

All of our computational artifacts can be built with the \texttt{carp-umbrella} metarepository (v1.7) located here:
\begin{description}[style=nextline, leftmargin=1.5cm, labelsep=0.0cm]
    \item[Zenodo:] https://doi.org/10.5281/zenodo.13314155
    \item[Github:] https://github.com/pdlfs/carp-umbrella
\end{description}

Our artifacts are:
\begin{description}
    \item[$A_1$] \texttt{vpic}
    \item[$A_2$] \texttt{particle.compressed.sample}
    \item[$A_3$] \texttt{range-runner + carp}
    \item[$A_4$] \texttt{compactor}
    \item[$A_5$] \texttt{range-reader}
\end{description}

\artexpl{
    Provide a table with the relevant computational artifacts,
    highlight their relation to the contributions (from above) and
    point to the elements in the paper that are reproducible by each artifact, e.g.,
    which figures or tables were generated with the artifact.
}

\artsampl{
    \begin{center}
        \begin{tabular}{rll}
            \toprule
            Artifact ID & Contributions & Related        \\
                        & Supported     & Paper Elements \\
            \midrule
            $A_1$       & $C_1$         & Tables 1-2     \\
                        &               & Figure 3       \\
            \midrule
            $A_2$       & $C_2$         & Tables 2-3     \\
                        &               & Figures 1-2    \\
            \midrule
            ..                                           \\
            \bottomrule
        \end{tabular}
    \end{center}
}

\begin{center}
    \begin{tabular}{rll}
        \toprule
        Artifact ID & Contributions & Related         \\
                    & Supported     & Paper Elements  \\
        \midrule
        $A_1$       & $C_2$         & Figure 1(a)     \\
        \midrule
        $A_2$       & $C_2$         & Figures 1(a), 8 \\
        \midrule
        $A_3$       & $C_1$         & Figure 7(a)     \\
        \midrule
        $A_4$       & $C_1$         & Figure 7(b)     \\
        \midrule
        $A_5$       & $C_1$         & Figure 7(b)     \\
        \bottomrule
    \end{tabular}
\end{center}


%%%%%%%%%%%%%%%%%%%%%%%%%%%%%%%%%%%%%%%%%%%%%%%%%%%%%%%%
\section{Artifact Identification}
%%%%%%%%%%%%%%%%%%%%%%%%%%%%%%%%%%%%%%%%%%%%%%%%%%%%%%%%
Running our artifacts at the scale described in the paper requires
an HPC cluster with significant compute, network, and storage resources.
However, it is possible to reduce the size of the simulation dataset
down to the point where smaller versions of our artifacts can be
replicated on a single node, as explained below.

% ----------------------------------------------
\newartifact

\artrel

VPIC an an open source MPI-based particle-in-a-cell simulation framework that
we used to generate the datasets indexed in the paper.  We provide
the two VPIC decks we used for our experiments: \texttt{reconnection}
and \texttt{reconnection512}. \texttt{reconnection} is a small 32-rank deck, while \texttt{reconnection512} requires 512 MPI ranks and generates the VPIC trace described in \S3.

\textbf{Running the reconnection decks is optional.}  Since running these decks is both time and resource intensive, we also provide a smaller micro-trace that can be used to run all experiments and demonstrate the functionality of CARP artifacts on a single node.

\artexp

The VPIC simulations create a \texttt{particle} directory and
periodically write particle output to it.  There is one output
file per rank for each timestep written.  For example, the
512 rank deck runs for 19,500 timesteps and writes particle
data every 200 timesteps.  Each timestep has 3.5B particles
(\textasciitilde200 GB).

\begin{example}
	particle
	|-- T.200
	|   |-- eparticle.200.0
	|   |-- eparticle.200.1
	|   |-- ...
	|   `-- eparticle.200.511
	|-- T.400
	|   `-- ...
	|-- T.600
	|   `-- ...
	`-- ...
\end{example}


\arttime

VPIC deck run times vary depending on both simulation parameters/state and the hardware configuration of the cluster used.   Expected run time from our cluster are:
\begin{itemize}
  \item \texttt{reconnection:} \textasciitilde600 minutes.  (The run can be terminated after \textasciitilde1 minute if only functional assessment is required).
  \item \texttt{reconnection512:} 1800-3600 minutes.
\end{itemize}

\artin

\artinpart{Hardware} A compute cluster with 512 cores across any number of nodes and a shared storage cluster with ~10TB of storage for \texttt{reconnection512}. A single node with 32 threads for \texttt{reconnection}.

\artinpart{Software} A standard Linux environment. While the codebase should build on most POSIX environments, we assume Ubuntu 20.04 LTS for the build tools described here. CARP also requires an MPI library --- while the code is compatible with all standard MPI implementations, the demo scripts use the \texttt{mpich}/\texttt{mvapich} command line interface (the demo scripts are not compatible with the \texttt{OpenMPI} CLI). Ensure that \texttt{mpirun}, \texttt{mpicc}, and \texttt{mpic++} are available via \texttt{\$PATH}.

\artinpart{Building}
All artifacts share a common build process that we will now explain.
Everything is built using a CMake-based metarepository called
\texttt{carp-umbrella}.  A \texttt{carp-umbrella} build pulls in
all CARP repositories (and their dependencies) and installs them
in the specified directory.  The following listing shows how to build
\texttt{carp-umbrella}.

\begin{bash}
# ensure mpich/mvapich are installed (not OpenMPI)
# ensure mpicc/mpirun in $PATH

# basic config + build tools follows:
/usr/sbin/sysctl kernel.yama.ptrace_scope=0
sudo apt-get install gcc g++ make cmake autoconf \
    automake libtool pkg-config git zip
    
# all artifacts will be installed here
INSTALL_TREE=$HOME/carp
mkdir -p $INSTALL_TREE/src
cd $INSTALL_TREE/src

# if using Zenodo zip:
unzip /path/to/carp-umbrella.zip \
    -d $INSTALL_TREE/src
    
# if using git repo:
git clone https://github.com/pdlfs/carp-umbrella

# by default, build will clone from repos
# see instructions in AE for self-contained build

mkdir carp-umbrella-build
cd carp-umbrella-build

# if using Zenodo zip, replace "../carp-umbrella" 
# with"../pdlfs-carp-umbrella-xxxxxxx" 
# for next command
cmake -DCMAKE_INSTALL_PREFIX=$INSTALL_TREE \
    -DMERCURY_OFI=OFF \
    -DMERCURY_POST_LIMIT=OFF \
    -DMERCURY_CHECKSUM=OFF \
    -DVPIC407=OFF \
    ../carp-umbrella
make -j
\end{bash}

The layout of the install tree is shown below.  Binaries are placed in \texttt{bin/} while scripts to interact with the artifacts are in \texttt{scripts/}.

\begin{example}
$INSTALL_TREE
   -- bin
   -- decks
   -- include
   -- lib
   -- scripts
   -- share
   -- src
       -- carp-umbrella
       -- carp-umbrella-build
\end{example}

\artcomp

We provide two scripts in \texttt{scripts/} --- \texttt{run\_vpic\_small.sh} and \texttt{run\_vpic\_big.sh} to run the two decks. The deck binaries are installed at \texttt{decks/}.

\texttt{run\_vpic\_small.sh} may be run on a single node. It will write its the trace to \texttt{/tmp/carp-demo/vpic-trace-small}. For functional evaluation of the artifact, we recommend running it for \textasciitilde1 minute, and then terminating it. As we provide a trace, it is not required to run it to completion. \texttt{run\_vpic\_big.sh} is an illustrative script, and it will need to be adjusted to match the requirements of the cluster it is run on.

\artout

\texttt{reconnection512} emits the trace described in Fig. 1 (a) of the paper, with structure and properties similar to artifact $A_2$. \texttt{reconnection} emits a trace that is similar, but the number of timesteps, particles etc. will be different.

% ----------------------------------------------

\newartifact

\artrel

This trace is a minified version of the trace generated by artifact $A_1$. It enables artifacts $A_3$, $A_4$, and $A_5$ to be run on a single compute node, with no need for a shared storage cluster.

\artexp
N/A
\arttime
N/A
\artin
Available at \texttt{\$INSTALL\_TREE/share/traces} once \texttt{carp-umbrella} is installed.
\artcomp
N/A
\artout

The trace has the following structure:

\begin{enumerate}
    \item It has three timesteps, in directies named \texttt{T.200}, \texttt{T.2000}, and \texttt{T.3800}.
    \item Each directory has 32 files named \texttt{eparticle.N} where $N (\in [0, 31]$) is the rank that wrote that file.
    \item Each $eparticle.N$ file contains a list of particle energies, each represented as a 4B little-endian float.
\end{enumerate}


% ----------------------------------------------

\newartifact

\artrel

\texttt{range-runner + carp} is the primary CARP artifact. \texttt{range-runner} loads the VPIC trace and replays it to simulate VPIC I/O. \texttt{carp} is built as a library that is preloaded via \texttt{LD\_PRELOAD} to intercept application I/O and index it in-situ for range-queries. While separate, the two can be considered a single artifact for evaluation purposes.

\artexp

If run at a large scale, in a storage-bound environment (as described in the paper, Fig 7b), the tool should take the same time with or without CARP.

At single-node scales the results are not expected to align with the paper. However, CARP-partitioned output can be used to trigger range queries with subsequent artifacts.

\arttime

\textasciitilde2 mins with \texttt{particle.compressed.micro}. More with larger applications.

\artin

\artinpart{Hardware} A single compute node with 32 CPU cores.

\artinpart{Software} A standard Linux environment with gcc, g++, and cmake. The metarepository \texttt{carp-umbrella} will pull in all other dependencies.

\artinpart{Datasets / Inputs} The artifact is set up to use the trace described as artifact $A_2$.

\artinpart{Building} See artifact $A_1$.

\artcomp

A script \texttt{run\_carp\_demo.sh} is provided to run the artifact. It is preconfigured to use the sample trace described in artifact $A_2$ and default values for various CARP parameters (which may be changed).

\artout

Example output for the provided VPIC trace:

\begin{example}
/path/to/particle
|-- RDB-00000000.tbl
|-- RDB-00000001.tbl
|-- RDB-00000002.tbl
|-- ...
\end{example}

As configured in the provided scripts, the artifact writes one file per CARP rank to a directory called \texttt{particle}. Each file is generated by a single \texttt{KoiDB} instance using a custom on-disk format. This output may be used for range queries using the artifact $A_5$.

% ----------------------------------------------

\newartifact

\artrel

\texttt{compactor} is used to generate a fully sorted version of the simulation output. It works by merging CARP's partially sorted output files into a full sorted and clustered index.
This output is used to measure the query performance of a sorted index in Fig 7a.

\artexp

\texttt{compactor} generates a copy of the trace data in a structure similar to the $A_2$ trace and the CARP output from $A_3$, except that the output is a fully sorted and clustered index.

The time taken by this artifact is not a proxy for the performance of parallel/distributed sorts. It is just an intermediate artifact to create fully sorted layouts without having to set up a sorting cluster.

\arttime

\textasciitilde1 min at single-node scale

\artin

We describe the CLI interface of this artifact here, but provide scripts to run it for convenience. This artifact needs to be run separately for each epoch in the trace.

\begin{example}
/path/to/compactor \
    -i $carp_output/particle \
    -o $sorted_output/$epoch \
    -e $epoch
\end{example}

\artcomp

\artinpart{Building} See artifact $A_1$.

\artinpart{Executing} This artifact is run within the demo workflow provided in \texttt{run\_carp\_demo.sh}.

It is preconfigured to use the CARP-partitioned trace output from artifact $A_3$.

\artout

Its output is particle data in a format identical to CARP, but fully sorted. For logistical reasons, it creates a separate directory for each epoch. Example output for the provided trace:

\begin{example}
/path/to/particle.sorted
|-- 0
|   `-- RDB-00000000.tbl
|-- 1
|   |-- RDB-00000000.tbl
`-- 2
    `-- RDB-00000000.tbl

\end{example}

% ----------------------------------------------

\newartifact

\artrel

\texttt{range-reader} is used to measure the query performance of both CARP and the fully sorted index, as in Fig 7a.

\artexp

\texttt{range-reader} takes a key range, and queries CARP or the fully sorted index for all particles in that range. It measures time taken to retrieve the particles and sort the particles sequentially.

\texttt{range-reader} can also be configured to process a batch of queries submitted as a CSV file, but we do not configure it that way in the scripts.

\arttime

\textasciitilde1 min with the provided trace.

\artin

We describe the CLI interface of this artifact here, but provide scripts to run it for convenience. If run with the \texttt{-a} flag, it analyzes the provided path and generates basic statistics. If run with a query, it executes the query and computes matching objects. If run with a query batch file, it reads the entire batch of queries from a CSV (format: \texttt{epoch,query\_begin,query\_end}).

\begin{example}
$INSTALL_TREE/bin/range-reader \
    -i <carp_out|sorted_out> \
    -p $PARALLELISM -a
    
$INSTALL_TREE/bin/range-reader \
    -i <carp_out|sorted_out> \
    -p $PARALLELISM -q \
    -e $query_epoch \
    -x $query_begin -y $query_end
    
$INSTALL_TREE/bin/range-reader \
    -i <carp_out|sorted_out> \
    -p $PARALLELISM -q \
    -b $query_batch.csv
\end{example}

\artcomp

\artinpart{Building} See artifact $A_1$.
\artinpart{Executing} Run as part of the workflow provided in \texttt{run\_carp\_demo.sh}

\artout

For a single query or analysis, this script writes output to \texttt{stdout}. For a query batch, it prints aggregated stats to \texttt{stdout}, and writes detailed query logs to a file called \texttt{querylog.csv}.

Query performance for CARP and the fully sorted index can be compared using this log, as in Fig 7a.



%%%%%%%%%%%%%%%%%%%%%%%%%%%%%%%%%%%%%%%%%%%%%%%%%%%%%
%  AE Appendix
%%%%%%%%%%%%%%%%%%%%%%%%%%%%%%%%%%%%%%%%%%%%%%%%%%%%%
\newpage
\appendixAE

\arteval{2}

\artin

Here, we provide some specific setup instructions tailored for Chamelon Cloud. We recommend a 32-thread node (such as the Skylake nodes on \texttt{CHI@TACC}) with the \texttt{Ubuntu20.04} image. The build process is the same as described in AD for artifact $A_1$ (note the differences when using a Zenodo zip vs using the git repository).

The artifact supports two build types --- via remote repositories, or via tarballs included in the release for a self-contained build. By default, following the instructions specified in the AD appendix will lead to specific commits being pulled in from the respective Git/HTTP URLs for all artifact repositories and dependencies. To build using the tarballs, execute these steps after extracting or cloning \texttt{carp-umbrella}.

\begin{bash}
# replace carp-umbrella with 
# pdlfs-carp-umbrella-xxx if using Zenodo zip
cd $INSTALL_TREE/src/carp-umbrella

ls -l cache.0 # ensure tars are present
# copy them to cache/ to "activate"
cp cache.0/*.gz cache/

cd $INSTALL_TREE/src
# follow usual build instructions from 
# "mkdir carp-umbrella-build" onwards
\end{bash}

\artcomp

We provide two scripts in \texttt{\$INSTALL\_TREE/scripts}:

\begin{enumerate}
    \item \texttt{run\_carp\_demo.sh} --- This script runs an interactive and guided end-to-end workflow starting with \texttt{range-runner + carp}. The sample trace is first partitioned via CARP, and then analyzed + queried using \texttt{range-reader}. The script then generates a fully sorted layout using \texttt{compactor}, and compares a CARP query vs a sorted query against the \texttt{compactor} output.
    \item \texttt{run\_carp\_demo\_intvl\_suite.sh} --- This script demonstrates CARP running with different renegotiation intervals, replicating the experiment described in {\S}VII-C (Pt. 4), at a smaller scale. It runs the artifact \texttt{range-runner + carp} with interval values of 2500, 5000, and 10,000.
\end{enumerate}

These scripts execute CARP with 16 ranks to permit wider compatibility. As the trace consists of 32 ranks, the number of MPI ranks can be changed if more cores are available.

\artout

\texttt{range-runner + carp} generates a partitioned output at \texttt{/tmp/carp-demo/plfs/particle}. The normalized load standard deviation is shown at the end of the output. A value of 0.05 means that CARP-generated partitions have a 5\% load imbalance (which is fairly balanced).

\arteval{3}

\artin
See $A_2$.

\artcomp
Run via \texttt{run\_carp\_demo.sh}, see $A_2$.

\artout

\texttt{compactor} generates fully sorted SSTs at \texttt{/tmp/carp-demo/plfs/particle.sorted/<epoch>}.

\arteval{4}

\artin

See $A_2$.
\artcomp
Run via \texttt{run\_carp\_demo.sh}, see $A_2$.

\artout

When run in analysis mode, \texttt{range-reader} shows that CARP partition selectivity is \textasciitilde6\%, at different points in the keyspace. That means that queries containing those keys will read at least \textasciitilde6\% data.

When run against a query, \texttt{range-reader} retrieves the same number of keys as the sorted layout.

\end{document}
